
\section{Conclusion}
\label{sec:conclusion}

\looseness=-1 We showed that learned post-hoc correction of frozen flow-matching VLAs is harmful by default. Under a K-matched comparison, unconditional correction degrades all four LIBERO suites on SmolVLA (0.4B) and all four on $\pi_{0.5}$ (3.6B), with the largest degradation on Long where the CRM is out-of-distribution. The root cause is that the vast majority of steps are already confident; corrections on these steps inject noise. Calibrated gating via $K{=}2$ pairwise $L_2$ distances detects confident steps and gates the CRM, recovering nearly all degradation on SmolVLA. The trigger rate shows a positive trend with improvement across suites, providing a built-in diagnostic. Because $K{=}2$ batched inference matches $K{=}1$ latency, the uncertainty probe is computationally free. The key finding is that \emph{knowing when not to correct} matters more than correction strength.
Task-agnostic CRM architectures that generalize beyond training tasks would remove the per-task retraining bottleneck; combining temporal correction~\cite{a2c2_2025} with our initial-action gating could address complementary failure modes; and real-robot validation under distribution shift would establish practical deployment viability.

\paragraph{Limitations.}
\label{sec:limitations}
\looseness=-1 The finding applies to the tested MLP-based CRM (${\sim}$327K parameters); larger correction architectures (e.g., transformer-based residual modules) may exhibit different failure modes. The $K{=}2$ probe requires diverse stochastic samples; deterministic VLAs such as OpenVLA-OFT produce near-identical outputs, yielding no usable uncertainty signal (Appendix~\ref{sec:applicability}), restricting applicability to flow-matching or diffusion-based policies. The threshold $\tau$ is a per-suite fixed value tuned on validation rollouts rather than learned end-to-end; an adaptive $\tau$ selection mechanism could improve robustness across task distributions. The CRM does not generalize to unseen tasks (leave-task-out SR ${\sim}$2\%), requiring practitioners to collect task-specific rollout data before applying correction, which limits plug-and-play deployment. Recovery on $\pi_{0.5}$ is partial (Table~\ref{tab:pi05_multiseed}), and all experiments use LIBERO simulation; real-robot transfer under visual and dynamics distribution shift remains future work. More broadly, the current CRM architecture does not surpass the VLA-only baseline---adaptive gating recovers the degradation introduced by unconditional correction but does not yield net performance gains beyond noise on SmolVLA, and only partial recovery on $\pi_{0.5}$. The primary contribution is therefore diagnostic and methodological (identifying when correction is harmful and providing a safe deployment framework) rather than achieving absolute performance improvement. Future work on task-agnostic correction architectures may close this gap.
